\begin{claim}{D.2.42}{Proved}{The actual secondary class lies outside the proposed odd-operation test}
Retain the notation of D.2.40, and put
\[
 L=Z'=Y\langle7\rangle,\qquad
 L'=Y\langle10\rangle,\qquad V=\Omega^4L'=U\langle6\rangle.
 \tag{D.2.42a}
\]
Thus $L$ is six-connected, $L'$ nine-connected and $V$
five-connected.  Let $y_8\in H^8L$ be the mod-two reduction
of the integral order-four coefficient connecting class of
the fundamental $\mathbb Z/4$ class in degree seven.
In the established basis it is $D_8^*+E_8^*$.
There is a common normalization of the relative square cocycle
which gives a class $\Theta_{15}\in H^{15}L'$ with
\[
 \sigma^{*4}\Theta_{15}=\mu_{11}|_V.
 \tag{D.2.42b}
\]
Here $\mu_{11}$ is the paired functional of D.2.40 and all
unlabeled cohomology has coefficients $\FF_2$.
The class is nonzero and satisfies
\[
 \mu_{11}|_V\notin
 \operatorname{Sq}^1H^{10}V+
 \operatorname{Sq}^3H^8V+
 \operatorname{Sq}^5H^6V.
 \tag{D.2.42c}
\]
This refutes that particular proposed sufficient operation
identity.  It does not assert that a map from the finite source
$D$ of D.2.38 realizes a nonzero evaluation, nor does it decide
the existence of a full fourth-suspension extension.
\end{claim}
\begin{proof}
We give the actual relative-cocycle construction before computing
its transgression.  This order distinguishes the class from a
formal bar symbol or an unspecified lift in a spectral sequence.

\medskip\noindent
\textbf{The square and its universal cocycle.}
The low double-loop cover basis of D.2.22 is
\[
 H_5E=\FF_2\{u\},\qquad
 H_6E=\FF_2\{D,E_6\},\qquad
 H_7E=\FF_2\{S,R\}.
 \tag{D.2.42d}
\]
The corresponding classes in $H_*(\Omega L)$, denoted
$u_6,D_7,E_7,S_8,R_8$, have images in the actual James
tensor algebra $H_*(\Omega Y)$
\[
 aa,\quad ab+ba,\quad f,\quad bb+g+h,\quad ac+ca,
 \tag{D.2.42e}
\]
respectively.  The tensor letters have degrees
$|a|=3$, $|b|=4$, $|c|=5$, $|f|=7$, $|g|=|h|=8$.
The complete eight ordered products of total degree fourteen
are linearly independent: the following table gives a tensor
word distinguishing each successive product from the others.
\[
\begin{array}{c|cccccccc}
\text{product}&uS&Su&uR&Ru&D^2&DE&ED&E^2\\ \hline
\text{word}&aag&gaa&aaac&caaa&abab&abf&fab&ff
\end{array}
\tag{D.2.42f}
\]
In a minimal Adams--Hilton model of the six-connected space
$L$, the least generator degree is six.  A differential of
degree fourteen can therefore have only these quadratic
terms; cubic terms start in degree eighteen.  Their independent
nonzero homology images exclude any nonzero such differential.
The quadratic-differential/reduced-diagonal identity gives
\[
 H^7L\,H^9L=0,\qquad H^8L\,H^8L=0
       \quad\text{in }H^{16}L.
 \tag{D.2.42g}
\]
The model and its diagonal interpretation are the same exact
ones used in D.2.26; see also
\cite[\S3.1, p.~300]{FGHHalperinLemaire1995}.
In particular $y_8^2=0$.

Choose an actual singular cocycle representing $y_8$ and a
cochain $c_{15}$ with $dc_{15}=y_8\smile y_8$.
For the cover map $q:L'\to L$, choose $P_7$ with
$dP_7=q^*y_8$, possible by nine-connectivity.  Then
\[
 \Theta_{15}=[q^*c_{15}-P_7\smile q^*y_8]
                     \in H^{15}L'
 \tag{D.2.42h}
\]
is closed.  A different $P_7$ changes it by a boundary, because
the difference is closed and $q^*y_8$ is exact.  Changing
$c_{15}$ by a cocycle $\ell_{15}$ changes $\Theta_{15}$ by
$q^*\ell_{15}$.
The actual singular-cochain Mayer--Vietoris construction of
D.2.26 says that its pullback by a lift of $h:S\to L$ is the
relative square class of the cone of $h$ whenever
$H^7S=H^8S=0$.  No abstract cone algebra is substituted for
the space in this assertion.

We now identify the required one common normalization.  Use
the universal evaluation map
\[
 \varepsilon:\Sigma^4V\longrightarrow L'\xrightarrow{q}L.
 \tag{D.2.42i}
\]
Its source starts in dimension ten.  For
$z\in\ker(\varepsilon_*:H_{15}(\Sigma^4V)\to H_{15}L)$,
take its connecting cone cycle and cancel the linear part of
the relative Adams--Hilton differential.  The remaining
degree-fourteen term has exactly the eight quadratic coordinates
in \textup{(D.2.42f)}.  Terms involving an old and a new
generator start in degree $6+10=16$, and old cubic terms
start in degree eighteen.  The quadratic-diagonal identity and
the actual cone pairing therefore identify the relative square
value on $z$ with the sum of the $D_7^2$ and $E_7^2$
coefficients of its once-adjoint image.  The mixed $D_7E_7$
and $E_7D_7$ coefficients are equal by primitivity and cancel;
the $uS,uR$ terms pair trivially with $y_8\otimes y_8$.
This is the per-vector argument of D.2.26j, in its explicitly
verified degree range.

Under \textup{(D.2.42e)}, the two square coordinates become
$(ab+ba)^2$ and $f^2$.  They are precisely the independent
restricted-square coordinates obtained by one homology
suspension from $\xi_1k$ and $\xi_1f$ in the free double-loop
basis.  Ordinary Hall-bracket coordinates and the other
restricted squares cannot change either coefficient; this is
the free basis used in D.2.29, from
\cite[Chapter~III, Theorems~1.2--1.4 and~3.1--3.2]{I_CohenLadaMay1976}.
Thus on the displayed kernel the relative square functional
equals $\mu_{11}|_V$, after removing the four source
suspensions.  Their difference consequently factors through
$\varepsilon_*H_{15}(\Sigma^4V)$.  Extend that functional to
one $\ell_{15}\in H^{15}L$ and replace $c_{15}$ by
$c_{15}+\ell_{15}$.  This proves \textup{(D.2.42b)} with
one normalization for all maps into $V$.
The fourfold cohomology suspension is an isomorphism here:
$L'$ is nine-connected and $V$ five-connected, and the
degree-eleven range is below the first product degree twelve.

\medskip\noindent
\textbf{The shifted Postnikov base.}
The proof of D.2.40, with the lower dimensions shifted by four,
gives the space-level description
\[
 P_9L\simeq P\times K(\FF_2,9),\qquad
 P=\operatorname{fib}\bigl(i\operatorname{Sq}^2\rho_2:
              K(\mathbb Z/4,7)\to K(\mathbb Z/4,9)\bigr).
 \tag{D.2.42j}
\]
For precision, the actual lower groups are
$\mathbb Z/4\{\alpha\}$ in degree seven,
$\mathbb Z/4\{\rho\}$ in degree eight, with
$\alpha\eta_7=2\rho$, and $\mathbb Z/2\{\lambda\}$
in degree nine.  Serre's stable integral calculation again
gives $H_8(K(\mathbb Z/4,7);\mathbb Z)=0$ and $H_9=\mathbb Z/2$,
so the unique nonzero lower invariant is the displayed one.
The homology suspension $H_7E\to H_9L$ is an isomorphism
in this range.  Thus the nonzero Hurewicz value $R$ of
$\lambda$ remains nonzero and excludes a nonzero next
Postnikov invariant, by exactly the transgression argument
of D.2.40g.  This proves \textup{(D.2.42j)} without assuming
an infinite-loop splitting.

Use $x_7,y_8$ for the fundamental and higher coefficient classes
on $K(\mathbb Z/4,7)$, $r_8,s_9$ for those of the fiber
$K(\mathbb Z/4,8)$, and $t_9$ for the fundamental class of
the separate $K(\FF_2,9)$ factor.
There is an actual class $r_8$ on $P$ with
\[
 \operatorname{Sq}^1x_7=\operatorname{Sq}^1y_8=0,
 \qquad\operatorname{Sq}^1r_8=\operatorname{Sq}^2x_7.
 \tag{D.2.42k}
\]
The last identity follows from the actual coefficient-fiber
model D.2.40h, with the dimensions shifted.  The pullback of
$y_8$ to $L$ is the class used above: the integral Moore-four
coordinates of D.2.22 identify its value as $D_8^*+E_8^*$.

The universal transgressive simple system of Serre gives the
following complete nonzero differential list with fiber source
degree at most fifteen:
\[
\begin{array}{c|c|c}
\text{source}&\text{differential}&\text{image}\\ \hline
s_9&d_{10}&\operatorname{Sq}^3x_7\\
\operatorname{Sq}^2s_9&d_{12}&\operatorname{Sq}^5x_7\\
\operatorname{Sq}^4s_9&d_{14}&
             \operatorname{Sq}^5\operatorname{Sq}^2x_7\\
\operatorname{Sq}^6s_9&d_{16}&
             \operatorname{Sq}^6\operatorname{Sq}^3x_7\\
\operatorname{Sq}^4\operatorname{Sq}^2s_9&d_{16}&
             \operatorname{Sq}^7\operatorname{Sq}^2x_7
\end{array}
\tag{D.2.42l}
\]
Here the higher coefficient image is
$\operatorname{Sq}^1\operatorname{Sq}^2x_7
=\operatorname{Sq}^3x_7$, not $x_7^2$.
Adem and $\operatorname{Sq}^1x_7=0$ give the later entries;
for the last also use instability of $\operatorname{Sq}^9x_7$.
The other relevant $s$-generators,
$\operatorname{Sq}^3s_9$ and $\operatorname{Sq}^5s_9$,
have zero transgression and survive in degrees twelve and
fourteen.  All $r$-generators survive as actual operations of
$r_8$.  The five images in \textup{(D.2.42l)} are independent
polynomial base generators and form a regular sequence.
Consequently no additional odd-annihilator cycles occur.
This uses the same precise simple-system theorem as D.2.40;
see \cite[\S2, no.~7, pp.~201--203, and no.~11,
pp.~204--205]{I_Serre1953}.

In particular the complete spanning lists for the three
degrees needed below are as follows; subscripts are suppressed
in the table.
\[
\begin{array}{c|l}
11&\operatorname{Sq}^4x,\ 
    \operatorname{Sq}^3y,\ \operatorname{Sq}^3r\\
13&\operatorname{Sq}^6x,\ \operatorname{Sq}^5y,\
    \operatorname{Sq}^4\operatorname{Sq}^2x,\
    \operatorname{Sq}^5r\\
15&\operatorname{Sq}^7y,\
    \operatorname{Sq}^6\operatorname{Sq}^2x,\
    \operatorname{Sq}^5\operatorname{Sq}^2y,\ xy,\\
  &\operatorname{Sq}^7r,\
    \operatorname{Sq}^5\operatorname{Sq}^2r,\ xr
\end{array}
\tag{D.2.42m}
\]
These are lists for $H^*(P)$, with no unspecified fiber lift
among their entries.  The extra permanent $s$-generators are
even-dimensional and produce no additional class in these
odd degrees: the first possible product involving $s_9$ is
$x_7s_9$ in degree sixteen, and it is nonpermanent; its
square begins in degree eighteen.
The separate $K(\FF_2,9)$ factor adds, respectively,
\[
\begin{array}{c|l}
11&\operatorname{Sq}^2t\\
13&\operatorname{Sq}^4t,\ 
            \operatorname{Sq}^3\operatorname{Sq}^1t\\
15&\operatorname{Sq}^6t,\
            \operatorname{Sq}^5\operatorname{Sq}^1t,\
            \operatorname{Sq}^4\operatorname{Sq}^2t.
\end{array}
\tag{D.2.42n}
\]
Cross-factor products start in degree $7+9=16$.

Apply the three odd squares to these complete lists.
On the $P$ factor, $\operatorname{Sq}^1H^{15}$ has only
the possible surviving value $x\operatorname{Sq}^2x$;
the other possibly nonzero value
$\operatorname{Sq}^7\operatorname{Sq}^2x$ is killed in
\textup{(D.2.42l)}.
Both $\operatorname{Sq}^3H^{13}(P)$ and
$\operatorname{Sq}^5H^{11}(P)$ give zero or that same killed
class.  For example
$\operatorname{Sq}^3\operatorname{Sq}^5r
=\operatorname{Sq}^7\operatorname{Sq}^1r
=\operatorname{Sq}^7\operatorname{Sq}^2x$ and
$\operatorname{Sq}^1(xr)=x\operatorname{Sq}^2x$.
On the separate factor the only resulting values are
$\operatorname{Sq}^7t$ and
$\operatorname{Sq}^5\operatorname{Sq}^2t$.
Therefore
\[
\begin{split}
 &\operatorname{Sq}^1H^{15}(P_9L)
  +\operatorname{Sq}^3H^{13}(P_9L)
  +\operatorname{Sq}^5H^{11}(P_9L)\\
 &\hspace{1cm}\subseteq
 \langle x\operatorname{Sq}^2x,
         \operatorname{Sq}^7t,
         \operatorname{Sq}^5\operatorname{Sq}^2t\rangle.
\end{split}
\tag{D.2.42o}
\]
The class $y_8^2$ is nonzero and independent of this span.
It is a polynomial square in the base of the first factor;
none of the five independent transgressions in
\textup{(D.2.42l)} kills it or relates it to another base
monomial.  Those two degree-fifteen fiber generators are also
the only possible sources of a differential into the scalar
base line in degree sixteen.  Thus its nonzero associated
filtration class survives to $H^{16}(P_9L)$.

\medskip\noindent
\textbf{Transgression of the actual relative class.}
In the actual fibration $L'\to L\to P_9L$, a mixed term
has total degree at least $7+10=17$.  Every fiber class of
degree $j$, $10\le j\le15$, is therefore transgressive:
for $r<7$ its outgoing differential has zero base target,
and for $7\le r\le j$ the target fiber degree is at most
nine and is zero.  The first possible differential is
$d_{j+1}$ to the scalar base line, and no earlier mixed
source can hit that line in this range.  Hence there are
unambiguous transgressions
\[
 \tau:H^jL'\longrightarrow H^{j+1}(P_9L),
             \qquad 10\le j\le15,
 \tag{D.2.42p}
\]
commuting with the Steenrod operations whose source and target
degrees are in this range.  This is ordinary cohomology
transgression compatibility, not a later-page operation with
an unmeasured indeterminacy.

For the particular class constructed above,
\[
                         \tau\Theta_{15}=y_8^2.
 \tag{D.2.42q}
\]
Indeed use a strict fibration model and choose $y_8$ on $L$
as the actual pullback of the base cocycle.  Its restriction
to the fiber is literally zero, so $P_7=0$ in
\textup{(D.2.42h)}.  The cochain $c_{15}$ satisfies
$dc_{15}=p^*(y_8^2)$, and the definition of transgression
gives \textup{(D.2.42q)}.  Adding the common closed
$\ell_{15}$ used to normalize \textup{(D.2.42b)} changes
the fiber class by a total-space restriction, whose
transgression is zero.

If $\Theta_{15}$ belonged to
$\operatorname{Sq}^1H^{14}L'+
\operatorname{Sq}^3H^{12}L'+
\operatorname{Sq}^5H^{10}L'$, operation compatibility in
\textup{(D.2.42p)} would put $y_8^2$ in the left side of
\textup{(D.2.42o)}, contradicting its proved independence.
The fourfold cohomology-suspension isomorphism commutes with
these operations, so \textup{(D.2.42b)} proves
\textup{(D.2.42c)}, including nonvanishing.
The operation condition was only a sufficient test for maps
from $D$.  Its failure does not settle the further attaching-map
and integral constraints on actual relative differences.
\end{proof}
